% =====================================================================
% Notas: Apartado 1 (Presentación del sistema) y Apartado 2 (Arquitectura)
% Responsable y expositor: Apunte Pazmiño Washington German
% Se incluye en notas/notas.tex con
% % =====================================================================
% Notas: Apartado 1 (Presentación del sistema) y Apartado 2 (Arquitectura)
% Responsable y expositor: Apunte Pazmiño Washington German
% Se incluye en notas/notas.tex con
% % =====================================================================
% Notas: Apartado 1 (Presentación del sistema) y Apartado 2 (Arquitectura)
% Responsable y expositor: Apunte Pazmiño Washington German
% Se incluye en notas/notas.tex con
% % =====================================================================
% Notas: Apartado 1 (Presentación del sistema) y Apartado 2 (Arquitectura)
% Responsable y expositor: Apunte Pazmiño Washington German
% Se incluye en notas/notas.tex con
% \input{notas/apartados/presentacion_arquitectura}
% Claves del .bib: spotifyAbout2026, spotifyQ32026, vansteen2023ds,
%   strid2023fleet, khlebnikova2024dataplatform, spotifyIncident2025,
%   gustavsson2023fleet, soylemez2022microservices
% Asistencia de IA (Claude): borrador revisado contra las fuentes y
%   reescrito con mis palabras.
% =====================================================================

\subsection*{Apartado 1: Presentación del sistema (expone Washington Apunte, 1 min)}

Spotify sirve para descubrir, organizar y escuchar música, pódcasts y
audiolibros desde una sola app. Tiene más de 100 millones de canciones y
8 millones de títulos de pódcast, y está en 184 mercados
\cite{spotifyAbout2026}. Lo usan dos grupos: quienes escuchan gratis con
anuncios y quienes pagan Premium. En el tercer trimestre de 2026 llegó a
788 millones de usuarios activos al mes, 11\,\% más que un año antes. De
ellos, 306 millones son suscriptores Premium y 499 millones usan la
versión con anuncios \cite[p.~3]{spotifyQ32026}.

Con esa escala no basta un servidor, ni siquiera un centro de datos
pequeño. Por eso Spotify es un sistema distribuido y corre en la nube.
La nube, dicho simple, es usar la infraestructura, la plataforma o las
aplicaciones de un proveedor como un servicio: IaaS, PaaS o SaaS
\cite[p.~99]{vansteen2023ds}. Spotify trabaja con Google Cloud desde
2016 y llevó allí su plataforma con \emph{lift and shift}, es decir,
moviendo lo que ya tenía \cite{strid2023fleet}. Antes corría sus
pipelines de datos en un clúster Hadoop propio, el más grande de Europa
\cite{khlebnikova2024dataplatform}.

\subsection*{Apartado 2: Arquitectura (expone Washington Apunte, 2 min)}

\paragraph{Componentes.}
Recorro el diagrama de izquierda a derecha. El usuario pulsa
«reproducir» y la app manda la petición por internet; si no recibe
respuesta, la vuelve a intentar \cite{spotifyIncident2025}. La petición
entra a una región de nube. La primera pieza que la recibe es el
perímetro, hecho con Envoy Proxy, y desde ahí pasa a los servicios de
backend \cite{spotifyIncident2025}. Esos servicios son miles de
componentes en producción, organizados como microservicios
\cite{gustavsson2023fleet}, y cada squad se encarga de los suyos
\cite{strid2023fleet}. Aparte, la app genera eventos de uso. La
plataforma de datos los recoge en colas Pub/Sub y los procesa en
BigQuery, Flink o Dataflow. Son más de un billón de eventos al día y más
de 38\,000 pipelines programados \cite{khlebnikova2024dataplatform}.

\paragraph{La nube en la arquitectura.}
En el diagrama se ven cuatro ideas de nube. La primera son las regiones:
el informe de la caída de 2025 habla de Europa, EE.\,UU. y Asia-Pacífico.
La segunda son los contenedores: Kubernetes ejecuta las instancias del
perímetro, cada una con un límite de memoria. El 16 de abril de 2025
Envoy podía usar más memoria que ese límite; Kubernetes apagaba las
instancias y el ciclo se repetía. Asia-Pacífico siguió funcionando porque
a esa hora tenía menos tráfico \cite{spotifyIncident2025}. La tercera es
la diferencia entre IaaS y PaaS. El \emph{lift and shift} se parece a
IaaS, porque es llevar tu sistema a la infraestructura del proveedor
\cite{strid2023fleet}. La plataforma de datos, en cambio, se apoya en
servicios administrados como Pub/Sub y BigQuery, y eso ya es PaaS
\cite{khlebnikova2024dataplatform}. La cuarta es la gestión a escala: con
Fleet Management, el parche de Log4j llegó al 80\,\% de los servicios de
backend en 9 horas \cite{gustavsson2023fleet}.

\paragraph{Ejemplo propio: SREI.}
SREI es nuestro PFC: un sistema web para registrar los eventos
interculturales de la UTEQ. Tiene un frontend en Angular que llama a una
sola API REST en Spring Boot, con PostgreSQL y los archivos guardados en
la carpeta \texttt{/uploads} del mismo servidor. Es cliente-servidor en
capas, centralizado, y lo desplegamos en local para las pruebas. Nuestro
propio requisito QR-003 pide que funcione sin grandes fallos en los días
laborables, pero si ese servidor se apaga en pleno evento, se cae todo:
es un punto único de fallo. No por eso tocaría copiar los microservicios
de Spotify. Adoptarlos no es trivial y sus retos se deben evaluar según
los requisitos de cada proyecto \cite[p.~1]{soylemez2022microservices}.
Para el tamaño de SREI, yo empezaría por sacar los archivos a un
almacenamiento aparte y usar una base de datos administrada.

% Claves del .bib: spotifyAbout2026, spotifyQ32026, vansteen2023ds,
%   strid2023fleet, khlebnikova2024dataplatform, spotifyIncident2025,
%   gustavsson2023fleet, soylemez2022microservices
% Asistencia de IA (Claude): borrador revisado contra las fuentes y
%   reescrito con mis palabras.
% =====================================================================

\subsection*{Apartado 1: Presentación del sistema (expone Washington Apunte, 1 min)}

Spotify sirve para descubrir, organizar y escuchar música, pódcasts y
audiolibros desde una sola app. Tiene más de 100 millones de canciones y
8 millones de títulos de pódcast, y está en 184 mercados
\cite{spotifyAbout2026}. Lo usan dos grupos: quienes escuchan gratis con
anuncios y quienes pagan Premium. En el tercer trimestre de 2026 llegó a
788 millones de usuarios activos al mes, 11\,\% más que un año antes. De
ellos, 306 millones son suscriptores Premium y 499 millones usan la
versión con anuncios \cite[p.~3]{spotifyQ32026}.

Con esa escala no basta un servidor, ni siquiera un centro de datos
pequeño. Por eso Spotify es un sistema distribuido y corre en la nube.
La nube, dicho simple, es usar la infraestructura, la plataforma o las
aplicaciones de un proveedor como un servicio: IaaS, PaaS o SaaS
\cite[p.~99]{vansteen2023ds}. Spotify trabaja con Google Cloud desde
2016 y llevó allí su plataforma con \emph{lift and shift}, es decir,
moviendo lo que ya tenía \cite{strid2023fleet}. Antes corría sus
pipelines de datos en un clúster Hadoop propio, el más grande de Europa
\cite{khlebnikova2024dataplatform}.

\subsection*{Apartado 2: Arquitectura (expone Washington Apunte, 2 min)}

\paragraph{Componentes.}
Recorro el diagrama de izquierda a derecha. El usuario pulsa
«reproducir» y la app manda la petición por internet; si no recibe
respuesta, la vuelve a intentar \cite{spotifyIncident2025}. La petición
entra a una región de nube. La primera pieza que la recibe es el
perímetro, hecho con Envoy Proxy, y desde ahí pasa a los servicios de
backend \cite{spotifyIncident2025}. Esos servicios son miles de
componentes en producción, organizados como microservicios
\cite{gustavsson2023fleet}, y cada squad se encarga de los suyos
\cite{strid2023fleet}. Aparte, la app genera eventos de uso. La
plataforma de datos los recoge en colas Pub/Sub y los procesa en
BigQuery, Flink o Dataflow. Son más de un billón de eventos al día y más
de 38\,000 pipelines programados \cite{khlebnikova2024dataplatform}.

\paragraph{La nube en la arquitectura.}
En el diagrama se ven cuatro ideas de nube. La primera son las regiones:
el informe de la caída de 2025 habla de Europa, EE.\,UU. y Asia-Pacífico.
La segunda son los contenedores: Kubernetes ejecuta las instancias del
perímetro, cada una con un límite de memoria. El 16 de abril de 2025
Envoy podía usar más memoria que ese límite; Kubernetes apagaba las
instancias y el ciclo se repetía. Asia-Pacífico siguió funcionando porque
a esa hora tenía menos tráfico \cite{spotifyIncident2025}. La tercera es
la diferencia entre IaaS y PaaS. El \emph{lift and shift} se parece a
IaaS, porque es llevar tu sistema a la infraestructura del proveedor
\cite{strid2023fleet}. La plataforma de datos, en cambio, se apoya en
servicios administrados como Pub/Sub y BigQuery, y eso ya es PaaS
\cite{khlebnikova2024dataplatform}. La cuarta es la gestión a escala: con
Fleet Management, el parche de Log4j llegó al 80\,\% de los servicios de
backend en 9 horas \cite{gustavsson2023fleet}.

\paragraph{Ejemplo propio: SREI.}
SREI es nuestro PFC: un sistema web para registrar los eventos
interculturales de la UTEQ. Tiene un frontend en Angular que llama a una
sola API REST en Spring Boot, con PostgreSQL y los archivos guardados en
la carpeta \texttt{/uploads} del mismo servidor. Es cliente-servidor en
capas, centralizado, y lo desplegamos en local para las pruebas. Nuestro
propio requisito QR-003 pide que funcione sin grandes fallos en los días
laborables, pero si ese servidor se apaga en pleno evento, se cae todo:
es un punto único de fallo. No por eso tocaría copiar los microservicios
de Spotify. Adoptarlos no es trivial y sus retos se deben evaluar según
los requisitos de cada proyecto \cite[p.~1]{soylemez2022microservices}.
Para el tamaño de SREI, yo empezaría por sacar los archivos a un
almacenamiento aparte y usar una base de datos administrada.

% Claves del .bib: spotifyAbout2026, spotifyQ32026, vansteen2023ds,
%   strid2023fleet, khlebnikova2024dataplatform, spotifyIncident2025,
%   gustavsson2023fleet, soylemez2022microservices
% Asistencia de IA (Claude): borrador revisado contra las fuentes y
%   reescrito con mis palabras.
% =====================================================================

\subsection*{Apartado 1: Presentación del sistema (expone Washington Apunte, 1 min)}

Spotify sirve para descubrir, organizar y escuchar música, pódcasts y
audiolibros desde una sola app. Tiene más de 100 millones de canciones y
8 millones de títulos de pódcast, y está en 184 mercados
\cite{spotifyAbout2026}. Lo usan dos grupos: quienes escuchan gratis con
anuncios y quienes pagan Premium. En el tercer trimestre de 2026 llegó a
788 millones de usuarios activos al mes, 11\,\% más que un año antes. De
ellos, 306 millones son suscriptores Premium y 499 millones usan la
versión con anuncios \cite[p.~3]{spotifyQ32026}.

Con esa escala no basta un servidor, ni siquiera un centro de datos
pequeño. Por eso Spotify es un sistema distribuido y corre en la nube.
La nube, dicho simple, es usar la infraestructura, la plataforma o las
aplicaciones de un proveedor como un servicio: IaaS, PaaS o SaaS
\cite[p.~99]{vansteen2023ds}. Spotify trabaja con Google Cloud desde
2016 y llevó allí su plataforma con \emph{lift and shift}, es decir,
moviendo lo que ya tenía \cite{strid2023fleet}. Antes corría sus
pipelines de datos en un clúster Hadoop propio, el más grande de Europa
\cite{khlebnikova2024dataplatform}.

\subsection*{Apartado 2: Arquitectura (expone Washington Apunte, 2 min)}

\paragraph{Componentes.}
Recorro el diagrama de izquierda a derecha. El usuario pulsa
«reproducir» y la app manda la petición por internet; si no recibe
respuesta, la vuelve a intentar \cite{spotifyIncident2025}. La petición
entra a una región de nube. La primera pieza que la recibe es el
perímetro, hecho con Envoy Proxy, y desde ahí pasa a los servicios de
backend \cite{spotifyIncident2025}. Esos servicios son miles de
componentes en producción, organizados como microservicios
\cite{gustavsson2023fleet}, y cada squad se encarga de los suyos
\cite{strid2023fleet}. Aparte, la app genera eventos de uso. La
plataforma de datos los recoge en colas Pub/Sub y los procesa en
BigQuery, Flink o Dataflow. Son más de un billón de eventos al día y más
de 38\,000 pipelines programados \cite{khlebnikova2024dataplatform}.

\paragraph{La nube en la arquitectura.}
En el diagrama se ven cuatro ideas de nube. La primera son las regiones:
el informe de la caída de 2025 habla de Europa, EE.\,UU. y Asia-Pacífico.
La segunda son los contenedores: Kubernetes ejecuta las instancias del
perímetro, cada una con un límite de memoria. El 16 de abril de 2025
Envoy podía usar más memoria que ese límite; Kubernetes apagaba las
instancias y el ciclo se repetía. Asia-Pacífico siguió funcionando porque
a esa hora tenía menos tráfico \cite{spotifyIncident2025}. La tercera es
la diferencia entre IaaS y PaaS. El \emph{lift and shift} se parece a
IaaS, porque es llevar tu sistema a la infraestructura del proveedor
\cite{strid2023fleet}. La plataforma de datos, en cambio, se apoya en
servicios administrados como Pub/Sub y BigQuery, y eso ya es PaaS
\cite{khlebnikova2024dataplatform}. La cuarta es la gestión a escala: con
Fleet Management, el parche de Log4j llegó al 80\,\% de los servicios de
backend en 9 horas \cite{gustavsson2023fleet}.

\paragraph{Ejemplo propio: SREI.}
SREI es nuestro PFC: un sistema web para registrar los eventos
interculturales de la UTEQ. Tiene un frontend en Angular que llama a una
sola API REST en Spring Boot, con PostgreSQL y los archivos guardados en
la carpeta \texttt{/uploads} del mismo servidor. Es cliente-servidor en
capas, centralizado, y lo desplegamos en local para las pruebas. Nuestro
propio requisito QR-003 pide que funcione sin grandes fallos en los días
laborables, pero si ese servidor se apaga en pleno evento, se cae todo:
es un punto único de fallo. No por eso tocaría copiar los microservicios
de Spotify. Adoptarlos no es trivial y sus retos se deben evaluar según
los requisitos de cada proyecto \cite[p.~1]{soylemez2022microservices}.
Para el tamaño de SREI, yo empezaría por sacar los archivos a un
almacenamiento aparte y usar una base de datos administrada.

% Claves del .bib: spotifyAbout2026, spotifyQ32026, vansteen2023ds,
%   strid2023fleet, khlebnikova2024dataplatform, spotifyIncident2025,
%   gustavsson2023fleet, soylemez2022microservices
% Asistencia de IA (Claude): borrador revisado contra las fuentes y
%   reescrito con mis palabras.
% =====================================================================

\subsection*{Apartado 1: Presentación del sistema (expone Washington Apunte, 1 min)}

Spotify sirve para descubrir, organizar y escuchar música, pódcasts y
audiolibros desde una sola app. Tiene más de 100 millones de canciones y
8 millones de títulos de pódcast, y está en 184 mercados
\cite{spotifyAbout2026}. Lo usan dos grupos: quienes escuchan gratis con
anuncios y quienes pagan Premium. En el tercer trimestre de 2026 llegó a
788 millones de usuarios activos al mes, 11\,\% más que un año antes. De
ellos, 306 millones son suscriptores Premium y 499 millones usan la
versión con anuncios \cite[p.~3]{spotifyQ32026}.

Con esa escala no basta un servidor, ni siquiera un centro de datos
pequeño. Por eso Spotify es un sistema distribuido y corre en la nube.
La nube, dicho simple, es usar la infraestructura, la plataforma o las
aplicaciones de un proveedor como un servicio: IaaS, PaaS o SaaS
\cite[p.~99]{vansteen2023ds}. Spotify trabaja con Google Cloud desde
2016 y llevó allí su plataforma con \emph{lift and shift}, es decir,
moviendo lo que ya tenía \cite{strid2023fleet}. Antes corría sus
pipelines de datos en un clúster Hadoop propio, el más grande de Europa
\cite{khlebnikova2024dataplatform}.

\subsection*{Apartado 2: Arquitectura (expone Washington Apunte, 2 min)}

\paragraph{Componentes.}
Recorro el diagrama de izquierda a derecha. El usuario pulsa
«reproducir» y la app manda la petición por internet; si no recibe
respuesta, la vuelve a intentar \cite{spotifyIncident2025}. La petición
entra a una región de nube. La primera pieza que la recibe es el
perímetro, hecho con Envoy Proxy, y desde ahí pasa a los servicios de
backend \cite{spotifyIncident2025}. Esos servicios son miles de
componentes en producción, organizados como microservicios
\cite{gustavsson2023fleet}, y cada squad se encarga de los suyos
\cite{strid2023fleet}. Aparte, la app genera eventos de uso. La
plataforma de datos los recoge en colas Pub/Sub y los procesa en
BigQuery, Flink o Dataflow. Son más de un billón de eventos al día y más
de 38\,000 pipelines programados \cite{khlebnikova2024dataplatform}.

\paragraph{La nube en la arquitectura.}
En el diagrama se ven cuatro ideas de nube. La primera son las regiones:
el informe de la caída de 2025 habla de Europa, EE.\,UU. y Asia-Pacífico.
La segunda son los contenedores: Kubernetes ejecuta las instancias del
perímetro, cada una con un límite de memoria. El 16 de abril de 2025
Envoy podía usar más memoria que ese límite; Kubernetes apagaba las
instancias y el ciclo se repetía. Asia-Pacífico siguió funcionando porque
a esa hora tenía menos tráfico \cite{spotifyIncident2025}. La tercera es
la diferencia entre IaaS y PaaS. El \emph{lift and shift} se parece a
IaaS, porque es llevar tu sistema a la infraestructura del proveedor
\cite{strid2023fleet}. La plataforma de datos, en cambio, se apoya en
servicios administrados como Pub/Sub y BigQuery, y eso ya es PaaS
\cite{khlebnikova2024dataplatform}. La cuarta es la gestión a escala: con
Fleet Management, el parche de Log4j llegó al 80\,\% de los servicios de
backend en 9 horas \cite{gustavsson2023fleet}.

\paragraph{Ejemplo propio: SREI.}
SREI es nuestro PFC: un sistema web para registrar los eventos
interculturales de la UTEQ. Tiene un frontend en Angular que llama a una
sola API REST en Spring Boot, con PostgreSQL y los archivos guardados en
la carpeta \texttt{/uploads} del mismo servidor. Es cliente-servidor en
capas, centralizado, y lo desplegamos en local para las pruebas. Nuestro
propio requisito QR-003 pide que funcione sin grandes fallos en los días
laborables, pero si ese servidor se apaga en pleno evento, se cae todo:
es un punto único de fallo. No por eso tocaría copiar los microservicios
de Spotify. Adoptarlos no es trivial y sus retos se deben evaluar según
los requisitos de cada proyecto \cite[p.~1]{soylemez2022microservices}.
Para el tamaño de SREI, yo empezaría por sacar los archivos a un
almacenamiento aparte y usar una base de datos administrada.
